\documentclass[aspectratio=169,11pt]{beamer}

% ============================================
% PACKAGES
% ============================================
\usepackage[utf8]{inputenc}
\usepackage[T1]{fontenc}
\usepackage[ngerman]{babel}
\usepackage{lmodern}
\usepackage{tcolorbox}
\tcbuselibrary{skins, hooks}
\usepackage{listings}
\usepackage{xcolor}
\usepackage{fontawesome5}
\usepackage{pgfplots}
\pgfplotsset{compat=1.18}
\usepackage{amsmath}

% Modern Font
\IfFileExists{FiraSans.sty}{\usepackage[sfdefault]{FiraSans}}{}
\renewcommand{\familydefault}{\sfdefault}

% ============================================
% COLORS & DESIGN SYSTEM
% ============================================
\definecolor{BgColor}{HTML}{FAFAFA}
\definecolor{Primary}{HTML}{2B3A42}
\definecolor{Accent}{HTML}{E84A5F}
\definecolor{Secondary}{HTML}{3F5765}
\definecolor{CodeBg}{HTML}{F0F0F0}
\definecolor{Success}{HTML}{28A745}
\definecolor{Warning}{HTML}{FFC107}

\setbeamercolor{background canvas}{bg=BgColor}
\setbeamercolor{title}{fg=Primary}
\setbeamercolor{frametitle}{bg=Primary, fg=white}
\setbeamercolor{structure}{fg=Accent}
\setbeamercolor{normal text}{fg=Primary}
\setbeamercolor{itemize item}{fg=Accent}

\setbeamertemplate{navigation symbols}{}
\setbeamertemplate{footline}{%
  \leavevmode%
  \hbox{%
  \begin{beamercolorbox}[wd=.333333\paperwidth,ht=2.25ex,dp=1ex,center]{author in head/foot}%
    \usebeamerfont{author in head/foot}\tiny Falko Himmel, Malte Baumann, Max Lehm
  \end{beamercolorbox}%
  \begin{beamercolorbox}[wd=.333333\paperwidth,ht=2.25ex,dp=1ex,center]{title in head/foot}%
    \usebeamerfont{title in head/foot}\tiny Grundlagen der Praktischen Informatik
  \end{beamercolorbox}%
  \begin{beamercolorbox}[wd=.333333\paperwidth,ht=2.25ex,dp=1ex,right]{date in head/foot}%
    \usebeamerfont{date in head/foot}\insertframenumber{} / \inserttotalframenumber\hspace*{2ex} 
  \end{beamercolorbox}}%
  \vskip0pt%
}

% ============================================
% CUSTOM BOXES
% ============================================
\tcbset{
    modernbox/.style={
        enhanced,
        boxrule=0pt,
        frame hidden,
        colback=white,
        coltitle=white,
        colbacktitle=Secondary,
        fonttitle=\bfseries,
        drop lifted shadow,
        arc=4pt,
        toptitle=2mm,
        bottomtitle=2mm,
        left=2mm,
        right=2mm
    },
    alertbox/.style={
        modernbox,
        colbacktitle=Accent
    },
    successbox/.style={
        modernbox,
        colbacktitle=Success
    }
}

% ============================================
% TITLE PAGE
% ============================================
\title{\textbf{Tutorium 10}}
\subtitle{Laufzeitkomplexität \& \texorpdfstring{$\mathcal{O}$}{O}-Notation}
\institute{Grundlagen der Praktischen Informatik}
\author{}
\date{\today}

\begin{document}

\begin{frame}[plain]
    \titlepage
\end{frame}


\begin{frame}{Inhaltsverzeichnis}
    \tableofcontents
\end{frame}


% ============================================
% THEORIE
% ============================================
\section{Laufzeitmessung: Praktisch vs. Theoretisch}

\begin{frame}{Wie messen wir Performance?}
    \begin{columns}[T]
        \begin{column}{0.48\textwidth}
            \begin{tcolorbox}[modernbox, title=Praktische Messung]
                Messung der tatsächlichen Zeit (Benchmarks).
                \begin{itemize}
                    \item[\textcolor{Success}{\faPlus}] Einfach umzusetzen
                    \item[\textcolor{Accent}{\faMinus}] Hardware- \& Umgebungsabhängig
                    \item[\textcolor{Accent}{\faMinus}] Schwer vergleichbar
                \end{itemize}
            \end{tcolorbox}
        \end{column}
        \begin{column}{0.48\textwidth}
            \begin{tcolorbox}[modernbox, title=Theoretische Analyse]
                Verhalten im asymptotischen Verlauf. Elementarschritte (Zuweisungen, Vergleiche, Add/Sub) zählen als konstante Zeit $\mathcal{O}(1)$.
                \begin{itemize}
                    \item[\textcolor{Success}{\faPlus}] Hardwareunabhängig
                    \item[\textcolor{Success}{\faPlus}] Fokus auf Skalierbarkeit
                \end{itemize}
            \end{tcolorbox}
        \end{column}
    \end{columns}
\end{frame}


\begin{frame}{Asymptotisches Verhalten}
    \begin{tcolorbox}[successbox, title=Wichtige Erkenntnisse für große $n$]
        \begin{enumerate}
            \item \textbf{Konstante Faktoren sind egal:} \\ Der Unterschied zwischen $2n$ und $10n$ ist im Vergleich zu $n^2$ vernachlässigbar.
            \item \textbf{Kleinere Terme verschwinden:} \\ In $n^2 + n + 5$ dominiert $n^2$ bei großen $n$ vollständig.
        \end{enumerate}
    \end{tcolorbox}
    
    \textbf{Rechenregeln:}
    \begin{itemize}
        \item $\mathcal{O}(c \cdot f(n)) = \mathcal{O}(f(n))$
        \item $\mathcal{O}(f(n) + g(n)) = \mathcal{O}(\max(f(n), g(n)))$
        \item $\mathcal{O}(f(n)) \cdot \mathcal{O}(g(n)) = \mathcal{O}(f(n) \cdot g(n))$
    \end{itemize}
\end{frame}



\begin{frame}{Asymptotisches Verhalten (Visuell)}
    \begin{columns}[T]
        \begin{column}{0.48\textwidth}
            \textbf{Konstante Faktoren sind egal}
            \begin{center}
            \begin{tikzpicture}[scale=0.55]
                \begin{axis}[
                    domain=0:10,
                    samples=50,
                    axis lines=left,
                    xlabel=$n$,
                    ylabel=$f(n)$,
                    legend pos=north west
                ]
                \addplot[thick, Accent] {x^2};
                \addlegendentry{$n^2$}
                \addplot[thick, Secondary] {2*x^2};
                \addlegendentry{$2n^2$}
                \addplot[thick, dashed, gray] {x^3};
                \addlegendentry{$n^3$}
                \end{axis}
            \end{tikzpicture}
            \end{center}
            \href{https://www.desmos.com/calculator/obdm27ccog}{\beamerbutton{In Desmos ansehen}}
        \end{column}
        \begin{column}{0.48\textwidth}
            \textbf{Kleine Terme verschwinden}
            \begin{center}
            \begin{tikzpicture}[scale=0.55]
                \begin{axis}[
                    domain=0:100,
                    samples=50,
                    axis lines=left,
                    xlabel=$n$,
                    legend pos=north west
                ]
                \addplot[thick, Accent] {x^2};
                \addlegendentry{$n^2$}
                \addplot[thick, Secondary] {x^2 + 10*x};
                \addlegendentry{$n^2 + 10n$}
                \end{axis}
            \end{tikzpicture}
            \end{center}
            \href{https://www.desmos.com/calculator/ok8ycuugjl}{\beamerbutton{In Desmos ansehen}}
        \end{column}
    \end{columns}
\end{frame}

\begin{frame}{Komplexitätsklassen}
    Typische Laufzeitklassen, geordnet nach Wachstum (am schnellsten zu am langsamsten):
    
    \vspace{0.3cm}
    \begin{tabular}{ll}
        \toprule
        \textbf{Klasse} & \textbf{Notation} \\
        \midrule
        Konstant & $\mathcal{O}(1)$ \\
        Logarithmisch & $\mathcal{O}(\log n)$ \\
        Linear & $\mathcal{O}(n)$ \\
        Log-Linear & $\mathcal{O}(n \log n)$ \\
        Quadratisch & $\mathcal{O}(n^2)$ \\
        Kubisch & $\mathcal{O}(n^3)$ \\
        Exponentiell & $\mathcal{O}(2^n)$ \\
        \bottomrule
    \end{tabular}
\end{frame}


% ============================================
% LIVE DEMO
% ============================================
\section{Praxis}

\begin{frame}[plain]
    \begin{center}
        \Huge \textbf{Live Demo} \\
        \vspace{0.5cm}
        \large \faPen~ \texttt{Altklausuraufgabe 2024\_1 (Aufgabe 6)} \\
        \vspace{0.3cm}
        \normalsize Wir analysieren zusammen die Laufzeitkomplexität am konkreten Beispiel!
\vspace{0.5cm}\begin{center}\href{https://moodle.uni-ulm.de/pluginfile.php/1348606/mod_folder/content/0/2024_1_exam.pdf}{\beamerbutton{Altklausur 2024\_1 auf Moodle öffnen}}\end{center}
    \end{center}
\end{frame}


\end{document}
